%%=============================================================================
%% Voorwoord
%%=============================================================================

\chapter*{\IfLanguageName{dutch}{Woord vooraf}{Preface}}%
\label{ch:voorwoord}

Dit onderzoek naar automatische netwerkdetectie bij Citymesh is ontstaan uit mijn interesse in netwerkbeheer en de uitdagingen die moderne IT-infrastructuren met zich meebrengen. Tijdens mijn opleiding System \& Network Administrator heb ik geleerd hoe belangrijk accurate netwerkdocumentatie is, maar ook hoe snel deze kan verouderen in dynamische omgevingen. De mogelijkheid om dit onderzoek uit te voeren bij Citymesh, een toonaangevende Belgische telecomoperator met complexe draadloze infrastructuren, bood een unieke kans om een praktisch relevant probleem te onderzoeken en een concrete oplossing te evalueren.

Ik wil graag mijn oprechte dank betuigen aan de mensen die hebben bijgedragen aan dit onderzoek. Allereerst aan mijn promotor, Martijn Saelens, voor zijn begeleiding, waardevolle feedback en ondersteuning gedurende het hele traject. Zijn inzichten hebben me geholpen om het onderzoek scherper te formuleren en de juiste focus te behouden.

Daarnaast ben ik zeer dankbaar voor de samenwerking met mijn co-promotor, Tom Pirlet van Citymesh. Zijn expertise in netwerkbeheer en zijn bereidheid om tijd vrij te maken voor interviews en overleg hebben dit onderzoek mogelijk gemaakt. De toegang tot de Citymesh-infrastructuur en de praktische inzichten die hij heeft gedeeld, hebben de relevantie en kwaliteit van dit onderzoek aanzienlijk verhoogd.

Ook wil ik mijn familie en vrienden bedanken voor hun steun en begrip tijdens de intensieve periode van onderzoek en schrijven. Hun aanmoediging heeft me geholpen om gemotiveerd te blijven, ook tijdens de uitdagende momenten.

Tenslotte hoop ik dat dit onderzoek een waardevolle bijdrage levert aan het domein van automatische netwerkdetectie en dat de resultaten en aanbevelingen nuttig zijn voor netwerkbeheerders die te maken hebben met vergelijkbare uitdagingen in hun organisaties.

Laurens Lecocq\\
Gent, \today
